% bank_fintech_attention/drafts/en/manuscript.tex
% P2 manuscript skeleton — M0 G0 scaffolded per B-route plan v2.55 (2026-05-17)
% Compile: xelatex + biber + xelatex × 2

% papers_assets/_shared/preamble_en.tex
% Shared English LaTeX preamble for all 5 paper drafts.
% P4.1 (plan-r3 derivative-paper extension, 2026-05-15)
% Compile: xelatex + biber + xelatex × 2

\documentclass[12pt,a4paper]{article}

% --- Fonts & Encoding ---
\usepackage{fontspec}
\setmainfont{Times New Roman}[Ligatures=TeX]

% --- Math ---
\usepackage{amsmath, amssymb, amsthm}

% --- Tables & Figures ---
\usepackage{booktabs}
\usepackage{graphicx}
\usepackage{longtable}
\usepackage{rotating}
\usepackage{caption}
\captionsetup{font=small,labelfont=bf,skip=4pt}

% --- Layout ---
\usepackage[margin=1in]{geometry}
\usepackage{setspace}
\onehalfspacing
\usepackage{hyperref}
\hypersetup{colorlinks=true,linkcolor=blue,citecolor=blue,urlcolor=blue}

% --- Bibliography (biblatex + biber) ---
\usepackage[backend=biber,style=apa,natbib=true]{biblatex}

% --- Useful macros ---
\newcommand{\code}[1]{\texttt{#1}}
\newcommand{\NoteThis}[1]{\textbf{\textcolor{red}{[TODO: #1]}}}
\usepackage[dvipsnames]{xcolor}

% --- Theorem environments ---
\newtheorem{prop}{Proposition}
\newtheorem{hyp}{Hypothesis}


\addbibresource{refs.bib}

\title{Bank-Level Fintech Attention Intensity:\\A Firm-Level Exposure Approach}
\author{\NoteThis{Author names}}
\date{\today}

\begin{document}
\maketitle

\begin{abstract}
\NoteThis{Final abstract pending — target $\approx$ 200 English words.}
Using bank-year media coverage of Chinese fintech and a firm-level attention-intensity
construct (\code{attention\_index\_full}) analogous to Hassan et~al.\ (2019, \emph{QJE}) on
political risk exposure, we study how cumulative fintech-topic exposure shapes bank-level
fintech investment, M\&A, patent output, and IT-staff growth across 18 years (2007--2024)
covering 500 Chinese commercial banks. We exploit five exogenous policy shocks (余额宝
launch / P2P regulation / Ant IPO pause / two data-export rules) for staggered DiD
identification, and instrument firm-level attention with a leave-one-out peer mean within
bank-type strata (with Bartik shift-share as fallback). We find \NoteThis{key result}\@.
\end{abstract}

\section{Introduction}\label{sec:intro}
\NoteThis{Motivate the research question and contribution.}

This paper differs from concurrent text-in-banking work along three margins.
First, attention is measured without polarity sign (cumulative exposure intensity
rather than tone direction). Second, identification exploits five China-specific
fintech policy shocks at high resolution. Third, the mediator is observable
(\code{ln\_fintech\_mdna\_length}, MD\&A textual length on fintech) rather than latent.

\section{Literature Review}\label{sec:lit}
Three strands: (1)~firm-level text-based exposure measurement
\autocite{hassan2019firm,florackis2023cybersecurity};
(2)~fintech taxonomy and innovation \autocite{gomber2018fintech,goldstein2019fintech};
(3)~text-in-banking exemplar \autocite{engelberg2018short}.

\section{Methodology}\label{sec:method}
\subsection{Sample and Data}
The bank--year panel covers $N \times T \approx 3000\text{--}9000$ bank--years
across 2007--2024 (per spec.sample). All bank--year observations are processed
through the \code{news\_process} pipeline.

\subsection{Attention Measurement (Stage 14)}
\NoteThis{Describe the four attention measures: index\_full / index\_ctx / disp / pct.}

\subsection{Identification Strategy}
\NoteThis{Staggered DiD (Callaway--Sant'Anna 2021) with leave-one-out peer-attention IV;
Bartik shift-share fallback.}

\subsection{Embedding-Based Measures (Co-Primary Dual Track; F5)}
All embedding-derived features are computed under two co-primary
embedding models from independent vendors, reported side by side as
mutual robustness: Alibaba \code{Qwen3-Embedding-8B} (4{,}096-d,
accessed via the SiliconFlow API; canonical unsuffixed columns) and
OpenAI \code{text-embedding-3-large} (3{,}072-d; \code{\_lg}-suffixed
columns), under strict column parity (each feature exists in both
tracks or in neither). Each track enters its own just-identified
specification; the two are never pooled into one over-identified
system. Earlier pipeline vintages used a Google Gemini embedding in
the canonical slot; as of pipeline v3.15.0 (June 2026) the canonical
slot is Alibaba \code{Qwen3-Embedding-8B}, with all embedding-derived
features regenerated under a controlled full re-embedding and
embedding-space semantic validity gated by a Spearman
rank-correlation check ($\rho \geq 0.80$). The sentiment-annotation
LLM ensemble (F1) is unaffected by this change: embedding models and
sentiment LLMs are disjoint pipeline components.

\section{Data}\label{sec:data}
\NoteThis{Data sources, sample selection, descriptive statistics (Table 1).}

\section{Results}\label{sec:results}
\NoteThis{Table 2 (M1$\to$M6 progressive main regression); Table 3 (mediation);
Table 4 (heterogeneity 4 bank types $\times$ fintech wave); Table 5 (robustness battery).}

\section{Conclusion}\label{sec:conclusion}
\NoteThis{Summarize findings and policy implications.}

\printbibliography

\end{document}
